% !TeX program = pdflatex
%==============================================================================
%  Lista Teórica E3 --- NLP + Classificação (Aplicação)
%  O gabarito sai de "Lista teorica E3 - gabarito.tex".
%==============================================================================
\providecommand{\gabaritoopt}{}
\documentclass[11pt,a4paper]{article}
\usepackage[\gabaritoopt]{../../recursos/latex/estilo-lista}

\begin{document}
\cabecalholista{E3}{NLP $+$ Classificação (Aplicação)}

%------------------------------------------------------------------------------
\begin{exercicio}[do texto à matriz]
Considere três documentos:
\begin{center}\small
$d_1 = $ ``o gato dorme'', \qquad
$d_2 = $ ``o cachorro dorme'', \qquad
$d_3 = $ ``o gato corre e o cachorro corre''.
\end{center}
\begin{enumerate}[label=(\alph*),leftmargin=1.1cm]
  \item Monte o vocabulário em ordem alfabética e escreva a matriz
        documento--termo de \textbf{contagens} ($3\times V$).
  \item O \texttt{scikit-learn} calcula
        $\mathrm{idf}(t) = \ln\!\big(\tfrac{1+n}{1+\mathrm{df}(t)}\big) + 1$,
        onde $\mathrm{df}(t)$ é o número de documentos que contêm $t$. Calcule o
        idf de \emph{gato} e de \emph{o}.
  \item Calcule o vetor TF-IDF de $d_1$, normalizado em norma $\ell_2$.
        \emph{Dados:} $\ln(4/3) \approx 0{,}2877$.
  \item Por que o termo \emph{o} recebe o menor peso possível? E por que ainda
        assim ele não é zerado por essa fórmula?
\end{enumerate}
\end{exercicio}

\begin{solucao}
\textbf{(a)} Vocabulário: \emph{cachorro, corre, dorme, e, gato, o} --- $V=6$.
\begin{center}\small
\renewcommand{\arraystretch}{1.25}
\begin{tabular}{@{}lcccccc@{}}
\toprule
      & cachorro & corre & dorme & e & gato & o \\
\midrule
$d_1$ & $0$ & $0$ & $1$ & $0$ & $1$ & $1$ \\
$d_2$ & $1$ & $0$ & $1$ & $0$ & $0$ & $1$ \\
$d_3$ & $1$ & $2$ & $0$ & $1$ & $1$ & $2$ \\
\bottomrule
\end{tabular}
\end{center}
Repare no que se perdeu: a \textbf{ordem} das palavras. ``o gato corre'' e
``corre o gato'' dão exatamente a mesma linha. É o que o nome \emph{bag of
words} --- saco de palavras --- quer dizer.

\smallskip
\textbf{(b)} Com $n=3$:
\begin{itemize}[leftmargin=1.4cm]
  \item \emph{gato} aparece em $d_1$ e $d_3$, então $\mathrm{df}=2$ e
        $\mathrm{idf} = \ln\!\big(\tfrac{4}{3}\big)+1 \approx \mathbf{1{,}2877}$;
  \item \emph{o} aparece nos três, $\mathrm{df}=3$ e
        $\mathrm{idf} = \ln\!\big(\tfrac{4}{4}\big)+1 = \ln 1 + 1 = \mathbf{1}$.
\end{itemize}

\smallskip
\textbf{(c)} O vetor bruto de $d_1$ é
$\text{tf}\times\text{idf}$, componente a componente. Só três entradas são não
nulas:
\[
  \text{dorme}: 1\times 1{,}2877, \qquad
  \text{gato}: 1\times 1{,}2877, \qquad
  \text{o}: 1\times 1 .
\]
A norma é
$\sqrt{1{,}2877^2+1{,}2877^2+1^2} = \sqrt{1{,}6582+1{,}6582+1} = \sqrt{4{,}3164}
\approx 2{,}0776$, e o vetor normalizado fica
\[
  \Big(0,\ 0,\ \underbrace{0{,}6198}_{\text{dorme}},\ 0,\
   \underbrace{0{,}6198}_{\text{gato}},\ \underbrace{0{,}4813}_{\text{o}}\Big).
\]

\smallskip
\textbf{(d)} Porque \emph{o} aparece em \textbf{todos} os documentos, e um termo
que aparece em todo lugar não distingue nada --- é exatamente isso que o idf
mede. Com $\mathrm{df}(t)=n$, o argumento do logaritmo é $\tfrac{1+n}{1+n}=1$ e o
idf cai ao seu valor mínimo.

Mas ele não é \emph{zero}: é $1$, por causa do ``$+1$'' na fórmula. Sem esse
termo, todo vocábulo presente em todos os documentos seria multiplicado por zero
e sumiria da matriz. O ``$+1$'' é uma decisão de projeto do
\texttt{scikit-learn} (o parâmetro \texttt{smooth\_idf} controla o outro ``$+1$'',
o de dentro do logaritmo) para que nenhum termo seja descartado \emph{pela
fórmula} --- descartar é papel do \texttt{min\_df} e da lista de
\emph{stop words}, que são decisões explícitas de quem modela.
\end{solucao}

%------------------------------------------------------------------------------
\begin{exercicio}[naive Bayes multinomial e a suavização de Laplace]
Um filtro de spam foi treinado sobre um vocabulário de $V=5$ palavras, com as
contagens totais abaixo (soma sobre todas as mensagens de cada classe):
\begin{center}\small
\renewcommand{\arraystretch}{1.25}
\begin{tabular}{@{}lccccc|c@{}}
\toprule
        & grátis & clique & oferta & reunião & amanhã & total \\
\midrule
spam    & $8$ & $5$ & $4$ & $0$ & $1$ & $18$ \\
não-spam& $1$ & $0$ & $1$ & $6$ & $7$ & $15$ \\
\bottomrule
\end{tabular}
\end{center}
As prioris estimadas são $\Prob(\text{spam})=0{,}4$ e
$\Prob(\text{não-spam})=0{,}6$. O modelo multinomial estima
\[
  \Prob(w \mid c) \;=\; \frac{N_{wc} + \alpha}{N_c + \alpha V},
\]
com $\alpha=1$ (suavização de Laplace).
\begin{enumerate}[label=(\alph*),leftmargin=1.1cm]
  \item Classifique a mensagem ``clique oferta grátis'', comparando as
        log-posterioris das duas classes.
        \emph{Dados:} $\ln 0{,}4 = -0{,}9163$, $\ln 0{,}6 = -0{,}5108$.
  \item O que aconteceria com a log-posteriori de não-spam se $\alpha=0$ (sem
        suavização)? Isso é razoável?
  \item Por que se trabalha com \emph{logaritmos} das probabilidades, e não com
        as probabilidades diretamente?
  \item O modelo supõe que as palavras são independentes dentro de cada classe.
        Cite duas maneiras óbvias pelas quais isso é falso em texto. Por que o
        método funciona mesmo assim?
\end{enumerate}
\end{exercicio}

\begin{solucao}
\textbf{(a)} Com $\alpha=1$ e $V=5$, o denominador é $18+5=23$ no spam e
$15+5=20$ no não-spam.
\begin{center}\small
\renewcommand{\arraystretch}{1.3}
\begin{tabular}{@{}lcc@{}}
\toprule
        & $\Prob(w\mid \text{spam})$ & $\Prob(w\mid \text{não-spam})$ \\
\midrule
clique  & $(5+1)/23 = 0{,}2609$ & $(0+1)/20 = 0{,}0500$ \\
oferta  & $(4+1)/23 = 0{,}2174$ & $(1+1)/20 = 0{,}1000$ \\
grátis  & $(8+1)/23 = 0{,}3913$ & $(1+1)/20 = 0{,}1000$ \\
\bottomrule
\end{tabular}
\end{center}
Somando os logaritmos:
\[
  \ln\Prob(\text{spam}\mid \text{msg}) \propto
  -0{,}9163 - 1{,}3437 - 1{,}5261 - 0{,}9383 = \mathbf{-4{,}7244},
\]
\[
  \ln\Prob(\text{não-spam}\mid \text{msg}) \propto
  -0{,}5108 - 2{,}9957 - 2{,}3026 - 2{,}3026 = \mathbf{-8{,}1117}.
\]
Como $-4{,}7244 > -8{,}1117$, a mensagem é classificada como \textbf{spam}. (A
diferença de $3{,}39$ em log corresponde a uma razão de chances de
$e^{3{,}39}\approx 30$ a favor de spam.)

\smallskip
\textbf{(b)} Sem suavização, $\Prob(\text{clique}\mid\text{não-spam}) =
0/15 = 0$, e o produto inteiro zera:
$\Prob(\text{não-spam}\mid\text{msg}) = 0$ \emph{exatamente}, ou $-\infty$ em
log.

Não é razoável. Uma única palavra nunca vista naquela classe no treino
\textbf{veta} a classe inteira, por mais que todas as outras palavras apontem
para ela. E ``nunca vista no treino'' é banal em texto: o Exercício~3 mostra que
a maioria dos termos aparece em pouquíssimos documentos. A suavização substitui a
afirmação ``é impossível'' por ``é raro'', que é o que os dados de fato
sustentam.

\smallskip
\textbf{(c)} Por \textbf{estabilidade numérica}. A posteriori é um produto de
tantos fatores quanto palavras na mensagem, cada um bem menor que 1. Numa
mensagem de 100 palavras com probabilidades típicas de $10^{-4}$, o produto é da
ordem de $10^{-400}$ --- abaixo do menor \emph{float} de precisão dupla
representável ($\approx 10^{-308}$), e portanto arredondado para zero. As duas
classes dariam zero, e a comparação seria impossível.

Em log, o produto vira soma e os valores ficam na casa das centenas negativas,
perfeitamente representáveis. Como o logaritmo é estritamente crescente,
comparar log-posterioris é equivalente a comparar posterioris.

\smallskip
\textbf{(d)} Duas maneiras óbvias:
\begin{itemize}[leftmargin=1.4cm]
  \item \textbf{colocações}: ``São'' e ``Paulo'' aparecem juntos muito mais do
        que o produto das frequências individuais previria; o mesmo para ``cartão
        de crédito'';
  \item \textbf{tema}: numa mensagem sobre futebol, ver ``gol'' aumenta muito a
        chance de ver ``time'' --- as palavras de um mesmo campo semântico se
        atraem, dentro da própria classe.
\end{itemize}

Funciona mesmo assim pela razão do Exercício~4 da Lista Teórica~08: para
\textbf{classificar} basta acertar qual log-posteriori é maior, e contar a mesma
evidência duas vezes tende a deslocar as duas na mesma direção, preservando a
ordem. Some-se a isso a economia de parâmetros --- o multinomial estima $V$
probabilidades por classe, e modelar dependências exigiria $V^2$ ou mais, o que é
inviável com $V$ na casa dos milhares.

O preço é a \textbf{calibração}, e é o mesmo medido na Lista prática~09: as
probabilidades saem empurradas para $0$ e $1$. Para decidir spam ou não, tanto
faz; para dizer ``esta mensagem tem 73\% de chance de ser spam'', não serve.
\end{solucao}

%------------------------------------------------------------------------------
\begin{exercicio}[esparsidade, e o que fazer com ela]
Numa coleção de $5\,572$ mensagens de SMS, o \texttt{CountVectorizer} produz um
vocabulário de $8\,672$ termos, e a matriz resultante tem $73\,916$ entradas não
nulas.
\begin{enumerate}[label=(\alph*),leftmargin=1.1cm]
  \item Calcule a densidade da matriz. Se ela fosse guardada densamente, em
        \emph{float} de 8 bytes, quanto ocuparia? E num formato esparso que
        guarda, para cada não-zero, o valor (8 bytes) e o índice da coluna
        (4 bytes)?
  \item Por que a matriz é tão esparsa? A mediana de palavras distintas por
        mensagem é $11$; use isso para prever a densidade e compare com o valor
        do item (a).
  \item O que faz \texttt{CountVectorizer(min\_df=5)}? Que efeito isso tem sobre
        o tamanho do vocabulário, e por que costuma \emph{ajudar} o
        classificador?
  \item Cite um método do curso que \textbf{não} funcionaria bem sobre uma matriz
        esparsa desse tipo, e explique por quê.
\end{enumerate}
\end{exercicio}

\begin{solucao}
\textbf{(a)} A densidade é
$73\,916 / (5\,572 \times 8\,672) = \mathbf{0{,}153\%}$ --- menos de duas
entradas em cada mil.

Densa: $5\,572 \times 8\,672 \times 8$ bytes $\approx 3{,}87\times 10^{8}$ bytes
$\approx \mathbf{387}$ \textbf{MB}. Esparsa: $73\,916 \times 12$ bytes
$\approx \mathbf{0{,}89}$ \textbf{MB} (mais um vetor de ponteiros de linha,
desprezível). Uma redução de \textbf{436 vezes}.

É o que torna o problema tratável num computador comum --- e é por isso que todos
os vetorizadores do \texttt{scikit-learn} devolvem matrizes
\texttt{scipy.sparse}, não \texttt{numpy}.

\smallskip
\textbf{(b)} Porque cada mensagem usa uma fração ínfima do vocabulário. Com
mediana de $11$ palavras distintas, a linha típica tem $11$ entradas não nulas em
$8\,672$ colunas, o que prevê densidade $11/8\,672 \approx 0{,}127\%$ --- próximo
dos $0{,}153\%$ medidos. (A diferença vem da cauda: algumas mensagens são bem
mais longas que a mediana, e a densidade é governada pela \emph{média}, não pela
mediana.)

O vocabulário é a união do que \emph{todas} as mensagens usam, e cada uma usa
quase nada dele.

Isso é uma propriedade estrutural de texto, não um acidente deste conjunto: o
vocabulário cresce com o tamanho do corpus enquanto o comprimento de cada
documento não.

\smallskip
\textbf{(c)} \texttt{min\_df=5} descarta todo termo que apareça em menos de 5
documentos. O vocabulário encolhe muito --- neste corpus, de $7\,168$ termos (no
conjunto de treino) para $1\,418$ --- um quinto. A razão é que $53{,}9\%$ dos
termos aparecem em \emph{um único} documento: nomes próprios, erros de digitação,
números.

Costuma ajudar por duas razões. Primeiro, \textbf{variância}: um termo que
aparece em 2 documentos, ambos spam, recebe um coeficiente enorme e confiante
baseado em duas observações --- é ruído que o modelo trata como sinal. Segundo,
\textbf{custo}: menos colunas, menos parâmetros, ajuste mais rápido.

O \texttt{min\_df} é, portanto, um hiperparâmetro de \textbf{regularização}
disfarçado de parâmetro de pré-processamento --- e, como tal, pertence ao
\texttt{Pipeline} e à grade de busca da Aula~07, não a uma decisão tomada uma vez
no começo.

\smallskip
\textbf{(d)} O \textbf{KNN} (ou qualquer método baseado em distância euclidiana).
Duas razões que se somam:

\begin{itemize}[leftmargin=1.4cm]
  \item é a \textbf{maldição da dimensionalidade} da Aula~05 na sua forma mais
        extrema: com $d\approx 8\,700$, a taxa $n^{-2/(2+d)}$ é indistinguível de
        $n^{0}$, e todos os pontos ficam praticamente equidistantes;
  \item pior, com esparsidade alta, duas mensagens quaisquer quase não
        compartilham palavras, então a distância entre elas é essencialmente
        $\sqrt{\lVert x_i\rVert^2 + \lVert x_j\rVert^2}$ --- determinada pelo
        \emph{comprimento} das mensagens, não pelo conteúdo delas.
\end{itemize}

Os métodos que funcionam bem aqui são os \textbf{lineares} --- naive Bayes,
logística, SVM linear --- justamente porque um modelo linear em $8\,700$
dimensões esparsas tem, na prática, pouquíssimos parâmetros ativos por
observação, e porque a penalização $\ell_2$ ou $\ell_1$ controla a variância sem
precisar de nenhuma noção de vizinhança.
\end{solucao}

%------------------------------------------------------------------------------
\begin{exercicio}[$n$-gramas e a explosão do vocabulário]
O \texttt{CountVectorizer} aceita \texttt{ngram\_range=(1,2)}, que acrescenta ao
vocabulário todos os \textbf{pares de palavras consecutivas} além das palavras
isoladas.
\begin{enumerate}[label=(\alph*),leftmargin=1.1cm]
  \item Para os três documentos do Exercício~1, liste os bigramas de cada um.
        Quantos bigramas distintos existem ao todo?
  \item Que informação os bigramas recuperam, que o saco de palavras havia
        jogado fora? Dê um exemplo em que isso muda a classificação.
  \item Um corpus com $V$ palavras distintas tem no máximo $V^2$ bigramas
        possíveis, mas na prática o número observado é muito menor. Por quê?
  \item Acrescentar bigramas aumenta muito o número de colunas. Isso agrava o
        problema do Exercício~3(d)? E o risco de superajuste?
\end{enumerate}
\end{exercicio}

\begin{solucao}
\textbf{(a)}
\begin{itemize}[leftmargin=1.4cm]
  \item $d_1$: \emph{o gato}, \emph{gato dorme};
  \item $d_2$: \emph{o cachorro}, \emph{cachorro dorme};
  \item $d_3$: \emph{o gato}, \emph{gato corre}, \emph{corre e}, \emph{e o},
        \emph{o cachorro}, \emph{cachorro corre}.
\end{itemize}
Distintos: \emph{o gato}, \emph{gato dorme}, \emph{o cachorro},
\emph{cachorro dorme}, \emph{gato corre}, \emph{corre e}, \emph{e o},
\emph{cachorro corre} --- \textbf{8 bigramas}. Com \texttt{ngram\_range=(1,2)} o
vocabulário passaria de $6$ para $6+8=\mathbf{14}$ colunas, mais que o dobro,
num corpus de três frases.

\smallskip
\textbf{(b)} Os bigramas recuperam, parcialmente, a \textbf{ordem} das palavras
--- a informação que o saco de palavras descartou.

O exemplo canônico em classificação de texto é a \textbf{negação}. As frases
``não é bom'' e ``é bom'' têm sacos de palavras quase idênticos (a segunda é um
subconjunto da primeira), e um modelo unigrama vê ``bom'' nas duas e tende a
classificar as duas como positivas. Com bigramas, ``não é'' e ``é bom'' são
colunas distintas, e o modelo pode aprender que ``não é'' carrega peso negativo.

Em detecção de spam, o mesmo vale para expressões: ``clique aqui'' é muito mais
informativo que ``clique'' e ``aqui'' separados.

\smallskip
\textbf{(c)} Porque a linguagem é fortemente restrita: a esmagadora maioria dos
pares de palavras \textbf{nunca ocorre}. Sintaxe (``o'' não é seguido de verbo),
semântica (``verde'' não modifica ``segunda-feira'') e simples convenção de uso
eliminam quase todas as $V^2$ combinações.

Empiricamente, o número de bigramas distintos num corpus cresce
aproximadamente com o número de \emph{tokens}, e não com $V^2$: é a lei de Heaps,
e a razão de fundo é a lei de Zipf --- a distribuição de frequências das palavras
tem cauda muito pesada, e a maioria delas aparece pouquíssimas vezes, logo tem
pouquíssimos vizinhos possíveis observados.

\smallskip
\textbf{(d)} \textbf{Para o KNN, agrava.} Mais colunas com a mesma quantidade de
informação por documento significa esparsidade ainda maior e distâncias ainda
menos informativas.

\textbf{Para o superajuste, o quadro é mais sutil}, e a resposta é ``depende do
que se faz com as colunas''. Sim, o número de parâmetros cresce, e um bigrama que
aparece em duas mensagens de spam vira um preditor perfeito no treino e ruído
fora dele --- exatamente o mecanismo do Exercício~3(c). Mas:
\begin{itemize}[leftmargin=1.4cm]
  \item \texttt{min\_df} corta justamente esses bigramas raríssimos, que são a
        maioria esmagadora dos novos;
  \item a penalização ($\ell_2$ da logística, ou o $C$ da SVM) controla os
        coeficientes das colunas que sobram;
  \item modelos lineares em alta dimensão esparsa toleram bem colunas extras
        inúteis, desde que regularizados --- é a observação de Greenshtein e
        Ritov citada na Aula~02, em que $d$ aparece só dentro de um logaritmo.
\end{itemize}
Na prática, \texttt{ngram\_range} e \texttt{min\_df} são escolhidos \emph{juntos},
por validação cruzada, dentro do mesmo \texttt{Pipeline} --- e é comum que a
combinação vencedora inclua bigramas com um \texttt{min\_df} mais alto do que o
usado só com unigramas.
\end{solucao}


\end{document}
